\documentclass{article}
\usepackage[T1]{fontenc}
\usepackage{lmodern}
\usepackage{graphicx}
\usepackage{amsmath,amssymb}
\usepackage[a4paper,margin=2cm]{geometry}
\usepackage[absolute,overlay]{textpos}
\setlength{\TPHorizModule}{1mm}
\setlength{\TPVertModule}{1mm}
\usepackage[indonesian]{babel}
\usepackage{hyperref}
\setlength{\emergencystretch}{2em}
% unit: Halaman 20

\begin{document}

% === Halaman 20 (python/native) ===

/ Program enkripsi dan dekripsi pesan dengan Caesar Cipher dalam Bahasa C++ 

\#include <iostream>

\#include <string.h>

using namespace std;

void enkripsi()

\{

  string plainteks, cipherteks;

  int i, k;

  char c;

  cout << "Ketikkan pesan:";

  cin.ignore(); getline (cin, plainteks);

  cout << "Masukkan jumlah pergesaran (0-25): "; cin >> k;

  cipherteks = "";  // inisialisasi cipherteks dengan null string


  for (i=0; i < plainteks.length(); i++) \{

      c = plainteks[i];

      if (isalpha(c)) \{    //hanya memproses huruf alfabet saja

         c = toupper(c);   // ubah menjadi huruf kapital 

         c = c - 65;       // kodekan huruf ke angka 0 s/d 25    

         c = (c + k) \% 26; // enkripsi, geser sejauh k ke kanan

         c = c + 65;       // kodekan kembali ke huruf semula          

      \}      

      cipherteks = cipherteks + c;  // sambungkan ke cipherteks

  \}

  cout << "Cipherteks: "<<cipherteks<< endl;  // cetak cipherteks

\}

20

Program Caesar Cipher dalam Bahasa C++

\end{document}
